\documentclass[10pt]{scrartcl}
\usepackage[mdthm,colorsec,noauthor,noasy]{evan}
\usepackage{handout}
\DeclareMathOperator{\im}{range}
\title{MATH 110 --- Week 8}
\subtitle{Notes Chapters 14--15; Homework 8}
\author{}\date{}\ihead{}
\begin{document}\maketitle
\section{Successive kernels and generalized eigenvectors}
\begin{qbox}What replaces eigenvectors when they do not span? For $N=\left(\begin{smallmatrix}0&1\\0&0\end{smallmatrix}\right)$, $Ne_2=e_1\ne0$ but $N^2e_2=0$: the second power detects the missing direction.\end{qbox}
\begin{prob}\textbf{Tools from earlier weeks.}\begin{itemize}\ii Weeks 2--4: basis extension and rank-nullity give $\dim V=\dim\ker A+\dim\im A$; an operator in finite dimension is invertible exactly when its kernel is zero.\ii Weeks 5--6: kernels and ranges of polynomials in $T$ are invariant, since $T$ commutes with its powers. Polynomial evaluation satisfies $(pq)(T)=p(T)q(T)$ by distributing finite sums.\ii Week 7: the minimal polynomial $m_T$ divides every annihilating polynomial, by division with remainder; $\deg m_T\le\dim V$ and its roots are exactly the eigenvalues. Over $\mathbb C$ it splits into linear factors by the fundamental theorem of algebra (Week 6, using the explicitly cited extreme-value theorem).\ii Week 7: an eigenbasis exists exactly when $m_T$ splits into distinct factors. Matrix columns give basis-vector images; triangular eigenvalues are diagonal entries, by back substitution.\end{itemize}\end{prob}
\begin{proposition*}[Notes Proposition 69; LADR 8.1--8.2]For any operator $A$, $\{0\}=\ker A^0\subseteq\ker A\subseteq\ker A^2\subseteq\cdots$. If $\ker A^m=\ker A^{m+1}$, then every later kernel equals $\ker A^m$.\end{proposition*}
\begin{motive}\ii[Want.] Know when further powers stop providing new vectors.\ii[Know.] Applying $A$ to zero gives zero.\ii[Key step.] Apply the assumed equality to $Av$ and propagate it one power at a time.\end{motive}
\begin{prob}\begin{proof}\begin{steps}\ii If $A^kv=0$, then $A^{k+1}v=A0=0$, proving each inclusion.\ii If $A^{m+2}v=0$, then $Av\in\ker A^{m+1}=\ker A^m$, hence $A^{m+1}v=0$ and $v\in\ker A^m$.\ii Thus $\ker A^{m+2}=\ker A^m$; repeat the same argument inductively for every later power.\end{steps}\end{proof}\end{prob}
\begin{prob}\begin{definition*}[Notes Definition 37; LADR 8.8, 8.19]For $T$ on finite-dimensional $V$, define $G(\lambda,T)=\bigcup_{k\ge0}\ker(T-\lambda I)^k$. Its nonzero vectors are generalized eigenvectors for $\lambda$: they satisfy $(T-\lambda I)^kv=0$ for some positive integer $k$.\end{definition*}\end{prob}
\begin{proposition*}[Notes Proposition 70; LADR 8.3, 8.20]If $n=\dim V$, then $G(\lambda,T)=\ker(T-\lambda I)^n$, an invariant subspace. The kernels of any operator $A$ stabilize by exponent $n$.\end{proposition*}
\begin{motive}\ii[Want.] Replace an infinite union by one kernel.\ii[Know.] A strictly increasing chain of subspaces increases dimension each time.\ii[Key step.] There cannot be $n+1$ strict increases inside dimension $n$.\end{motive}
\begin{prob}\begin{proof}\begin{steps}\ii If $\ker A^n\ne\ker A^{n+1}$, Proposition 69 implies every inclusion from $\ker A^0$ to $\ker A^{n+1}$ is strict. The last kernel would have dimension at least $n+1$, impossible.\ii Thus $\ker A^n=\ker A^{n+1}$, and Proposition 69 gives equality with every later kernel. Earlier kernels are contained in it.\ii Use $A=T-\lambda I$ to identify the union. It is invariant because $T$ commutes with $(T-\lambda I)^n$ (Week 6). The zero-dimensional case gives $\ker I=V=\{0\}$.\end{steps}\end{proof}\end{prob}
\begin{proposition*}[LADR 8.4: stabilized kernel and range]For $n=\dim V$ and $A\in L(V)$, $V=\ker A^n\oplus\im A^n$. Both summands are invariant under every operator commuting with $A$.\end{proposition*}
\begin{motive}\ii[Want.] Separate the vectors that a high power sends to zero from a complementary invariant space.\ii[Know.] The kernels at exponents $n$ and $2n$ agree.\ii[Key step.] A vector in both spaces would be the image of a vector already in $\ker A^n$.\end{motive}
\begin{prob}\begin{proof}\begin{steps}\ii If $w=A^nv$ and $A^nw=0$, then $A^{2n}v=0$. Stabilization gives $A^nv=0$, hence $w=0$.\ii Rank-nullity gives $\dim\ker A^n+\dim\im A^n=n$. The zero intersection therefore makes their sum all of $V$, by Week 4's dimension formula.\ii If $B$ commutes with $A$, then $A^nBv=BA^nv$ preserves the kernel, and $BA^nv=A^nBv$ preserves the range.\end{steps}\end{proof}\end{prob}
\begin{example*}[Notes Example 61]For $A=\left(\begin{smallmatrix}1&1&1\\0&1&1\\0&0&3\end{smallmatrix}\right)$ on $\mathbb Q^3$:\begin{steps}\ii $A-I=\left(\begin{smallmatrix}0&1&1\\0&0&1\\0&0&2\end{smallmatrix}\right)$ has kernel $\mathbb Qe_1$.\ii $(A-I)^2=\left(\begin{smallmatrix}0&0&3\\0&0&2\\0&0&4\end{smallmatrix}\right)$ and $(A-I)^3=\left(\begin{smallmatrix}0&0&6\\0&0&4\\0&0&8\end{smallmatrix}\right)$ both have kernel $\operatorname{span}(e_1,e_2)$.\ii Proposition 69 gives $G(1,A)=\operatorname{span}(e_1,e_2)$; $e_2$ is generalized but not an ordinary eigenvector.\end{steps}\end{example*}
\begin{example*}[LADR 8.6, 8.10, 8.21]\begin{steps}\ii For $T(x_1,x_2,x_3)=(4x_2,0,5x_3)$ over $\mathbb R$ or $\mathbb C$, $\ker T=Fe_1$ and $\im T=\operatorname{span}(e_1,e_3)$, so their intersection is nonzero.\ii $T^3(x)=(0,0,125x_3)$ gives $\ker T^3=\operatorname{span}(e_1,e_2)$ and $\im T^3=Fe_3$, now complementary.\ii Solving the shifted-power kernels gives $G(0,T)=\operatorname{span}(e_1,e_2)$ and $G(5,T)=Fe_3$. Ordinary eigenspaces are $Fe_1$ and $Fe_3$, so they do not span; the generalized eigenspaces do.\end{steps}\end{example*}
\begin{proposition*}[Notes Proposition 71]$G(\lambda,T)\ne\{0\}$ exactly when $\lambda$ is an eigenvalue.\end{proposition*}
\begin{motive}\ii[Want.] Check that generalized vectors do not introduce new eigenvalues.\ii[Know.] Ordinary eigenvectors already belong to the union of kernels.\ii[Key step.] If the first kernel is zero, Proposition 69 makes every later kernel zero.\end{motive}
\begin{prob}\begin{proof}\begin{steps}\ii If $\lambda$ is not an eigenvalue, $\ker(T-\lambda I)=\{0\}=\ker(T-\lambda I)^0$, so stabilization gives $G(\lambda,T)=\{0\}$.\ii If it is an eigenvalue, its nonzero eigenspace is contained in $G(\lambda,T)$ by definition.\end{steps}\end{proof}\end{prob}
\begin{proposition*}[Notes Proposition 72; LADR 8.11--8.12]Nonzero generalized eigenvectors for distinct eigenvalues are independent. In particular a generalized eigenvector corresponds to a unique eigenvalue, and the sum of distinct generalized eigenspaces is direct.\end{proposition*}
\begin{motive}\ii[Want.] Separate contributions from different eigenvalues.\ii[Know.] A sufficiently high shifted power sends each matching generalized vector to zero, and all the shifted powers commute.\ii[Key step.] Apply a product which makes all but one term zero; turn the remaining term into an ordinary eigenvector first.\end{motive}
\begin{prob}\begin{proof}\begin{steps}\ii Suppose $\sum_{j=1}^ra_jv_j=0$ with $v_j\in G(\lambda_j,T)\setminus\{0\}$. Fix an index $i$ and choose the largest $h$ for which $w=(T-\lambda_iI)^hv_i\ne0$; it exists by the defining eventual zero property.\ii Then $Tw=\lambda_iw$. For $j\ne i$, $(T-\lambda_jI)^nw=(\lambda_i-\lambda_j)^nw$, with $n=\dim V$.\ii Apply $(T-\lambda_iI)^h\prod_{j\ne i}(T-\lambda_jI)^n$ to the relation. Commutation and Proposition 70 make every term except $i$ zero.\ii The remaining equation is $a_i\prod_{j\ne i}(\lambda_i-\lambda_j)^nw=0$. Each scalar factor and $w$ is nonzero, so $a_i=0$. Do this for every $i$.\ii If a nonzero $v$ corresponded to two distinct eigenvalues, the dependent pair $v,v$ would contradict independence. Grouping any relation with one vector from each generalized eigenspace proves directness.\end{steps}\end{proof}\end{prob}
\begin{shortcuts}\ii Stop taking powers as soon as two consecutive kernels agree, by Proposition 69.\ii Use exponent $\dim V$ as a guaranteed stopping point, by Proposition 70.\ii A generalized eigenvector can be checked by a power equation even when it is not an ordinary eigenvector.\end{shortcuts}
\begin{trapbox}\ii A union of arbitrary subspaces need not be a subspace: the union of the coordinate axes omits $(1,1)$. The nested-kernel union is a subspace because it equals one stabilized kernel.\ii $\ker A$ and $\im A$ need not be complementary: the two-by-two $N$ in the question has both equal to $Fe_1$. Use stabilized powers.\end{trapbox}
\section{Nilpotent operators and chains}
\begin{qbox}How can one choose coordinates for an operator whose power is zero? For differentiation on $\mathcal P_2(\mathbb R)$, the ordered list $1,x,x^2/2$ has images $0,1,x$, giving ones above the diagonal.\end{qbox}
\begin{prob}\begin{definition*}[Notes Definition 39; LADR 8.14]An operator $N$ is nilpotent if $N^r=0$ for some positive integer $r$. On a nonzero finite-dimensional space its nilpotency index is the smallest such $r$. A matrix is strictly upper triangular if it is upper triangular with every diagonal entry zero. A chain of length $r$ is an independent ordered list $v_1,\ldots,v_r$ satisfying $Nv_1=0$ and $Nv_j=v_{j-1}$ for $j\ge2$.\end{definition*}\end{prob}
\begin{proposition*}[LADR 8.16--8.18: nilpotence tests]For $n=\dim V>0$, a nilpotent $N$ satisfies $N^n=0$ and has only eigenvalue 0. It has minimal polynomial $x^r$, where $r$ is its index. Equivalently, $N$ has a strictly upper-triangular matrix in some basis. Over $\mathbb C$, having only eigenvalue 0 also implies nilpotence.\end{proposition*}
\begin{motive}\ii[Plan.] Stabilization bounds the exponent; divisibility identifies the polynomial; an invariant-kernel basis gives triangular form.\ii[Want.] Recognize nilpotence from powers, polynomial relations or a matrix.\ii[Know.] Kernels stabilize by dimension, and minimal polynomials divide every annihilating polynomial (Week 7).\ii[Key step.] Restrict to a one-dimensional kernel subspace and induct on the quotient.\end{motive}
\begin{prob}\begin{proof}\begin{steps}\ii If $N^r=0$, then $\ker N^r=V$; stabilization implies $\ker N^n=V$, hence $N^n=0$. If $Nv=\lambda v$ with $v\ne0$, $N^rv=\lambda^rv=0$ forces $\lambda=0$; a nonzero nilpotent space has nonzero kernel, so 0 occurs.\ii Since $m_N\mid x^r$, it is a monic power $x^s$. Minimality of the index gives $s=r$. Conversely $m_N=x^r$ immediately gives nilpotence.\ii For a triangular basis, induct on dimension: choose $v_1\ne0$ in $\ker N$. The quotient by $Fv_1$ is nilpotent, so a triangular basis there lifts to one here with all diagonal entries zero.\ii Conversely, a strictly upper-triangular matrix maps the first $j$ basis vectors into the first $j-1$ span; $n$ applications give zero.\ii Over $\mathbb C$, if the only eigenvalue is 0, complex factorization and the root characterization of $m_N$ (Week 7) give $m_N=x^r$. On the zero space use $m_N=1$, the empty basis, and $N=0$.\end{steps}\end{proof}\end{prob}
\begin{example*}[LADR 8.15]\begin{steps}\ii On $F^4$, $N(x_1,x_2,x_3,x_4)=(0,0,x_1,x_2)$ has $N^2=0$, so is nilpotent of index 2.\ii Over $\mathbb R$ or $\mathbb C$, $B=\left(\begin{smallmatrix}-3&9&0\\-7&9&6\\4&0&-6\end{smallmatrix}\right)$ satisfies $B^2=\left(\begin{smallmatrix}-54&54&54\\-18&18&18\\-36&36&36\end{smallmatrix}\right)\ne0$ and $B^3=0$, so has index 3.\ii Differentiation on $\mathcal P_m(\mathbb R)$ has index $m+1$: the $(m+1)$st derivative is zero and $D^mx^m=m!\ne0$.\end{steps}\end{example*}
\begin{proposition*}[Chain independence; LADR 8A Exercise 2]If $N^rv=0$ and $N^{r-1}v\ne0$, then $N^{r-1}v,N^{r-2}v,\ldots,v$ is a chain of length $r$.\end{proposition*}
\begin{motive}\ii[Want.] Obtain independent vectors by repeated application.\ii[Know.] The last nonzero iterate cannot be a nonzero scalar multiple of zero.\ii[Key step.] Apply a power that isolates the earliest nonzero coefficient of a proposed relation.\end{motive}
\begin{prob}\begin{proof}\begin{steps}\ii In a relation $\sum_{j=0}^{r-1}a_jN^jv=0$, suppose $a_h$ is the first nonzero coefficient. Apply $N^{r-1-h}$.\ii Terms with $j>h$ vanish and the remaining term is $a_hN^{r-1}v=0$, impossible. Thus all coefficients are zero; the chain equations follow from the powers.\end{steps}\end{proof}\end{prob}
\begin{proposition*}[Notes Proposition 75; LADR 8.45]Every nilpotent operator has a basis consisting of disjoint chains. Its matrix is zero except for possible ones immediately above the diagonal.\end{proposition*}
\begin{motive}\ii[Plan.] Build chains in the smaller range, lift their last vectors, then add the missing kernel directions.\ii[Want.] Produce enough chains to span the entire space.\ii[Know.] The range of a nilpotent operator is proper and invariant, and a basis of its kernel can be extended.\ii[First try.] Extending one chain to an arbitrary basis gives no control of the images of the added vectors, so it need not produce a basis of chains.\ii[Key step.] Lifting the last vector of each range chain extends that chain by one vector.\end{motive}
\begin{prob}\begin{proof}\begin{steps}\ii Induct on dimension, with the empty basis for dimension 0. If $N=0$, any basis consists of length-1 chains. Otherwise $U=\im N$ is proper, and $N|_U$ is nilpotent.\ii By induction choose chains $N^{r_j-1}v_j,\ldots,v_j$ forming a basis of $U$. Choose $u_j$ with $Nu_j=v_j$; their lifted lists are $N^{r_j}u_j,\ldots,u_j$.\ii These lists are independent: apply $N$ to a proposed relation. The old basis of $U$ forces all coefficients except those of $N^{r_j}u_j$ to zero; those remaining vectors are independent members of the old basis.\ii The vectors $N^{r_j}u_j=N^{r_j-1}v_j$ form a basis of $\ker N\cap U$, by applying $N$ to the old chain basis. Extend them to a basis of $\ker N$ by adding $w_1,\ldots,w_\ell$. The lifted lists plus the $w_i$ are independent by the same argument.\ii The image of their span under $N$ is all of $U$. For any $v$, choose $z$ in that span with $Nz=Nv$; then $v-z\in\ker N$, also in the span. Thus the lists span $V$ and give the asserted matrix.\end{steps}\end{proof}\end{prob}
\begin{fig}\begin{tikzpicture}[>=stealth]\node(a)at(0,0){$v_3$};\node(b)at(2,0){$v_2$};\node(c)at(4,0){$v_1$};\node(d)at(6,0){$0$};\draw[->](a)--node[above]{$N$}(b);\draw[->](b)--node[above]{$N$}(c);\draw[->](c)--node[above]{$N$}(d);\end{tikzpicture}\figcap{In the ordered basis $(v_1,v_2,v_3)$ this chain contributes a three-by-three matrix with ones just above the diagonal.}\end{fig}
\begin{example*}[Notes Proposition 75; LADR 8.42--8.43]\begin{steps}\ii On $\mathbb C^4$, $T(x)=(0,x_1,x_2,x_3)$ has chain
\[(T^3e_1,T^2e_1,Te_1,e_1)=(e_4,e_3,e_2,e_1).\]
Its matrix in that basis has three superdiagonal ones.\ii On $\mathbb C^6$, $T(x)=(0,x_1,x_2,0,x_4,0)$ has chains $(e_3,e_2,e_1)$, $(e_5,e_4)$ and $(e_6)$. Their lengths are 3, 2 and 1.\ii The two four-dimensional matrices displayed with Notes Proposition 75 have chains of lengths $(2,2)$ and $(4)$, respectively: their nonzero superdiagonal entries are in positions $(1,2),(3,4)$ and $(1,2),(2,3),(3,4)$. Their indices are 2 and 4.\end{steps}\end{example*}
\begin{shortcuts}\ii A longest chain has length equal to the nilpotency index, by its last nonzero iterate and the chain-basis theorem.\ii To turn derivatives into unit chain steps, use $1,x,x^2/2!,\ldots,x^m/m!$.\end{shortcuts}
\begin{trapbox}\ii Over $\mathbb R$, only real eigenvalue 0 does not imply nilpotence: $\left(\begin{smallmatrix}0&0&0\\0&0&-1\\0&1&0\end{smallmatrix}\right)$ has only real eigenvalue 0 but its square is nonzero and its rotation part never becomes zero.\ii Nilpotent does not mean zero: $\left(\begin{smallmatrix}0&1\\0&0\end{smallmatrix}\right)$ is a nonzero example.\end{trapbox}
\section{Generalized eigenspace decomposition and characteristic polynomial}
\begin{qbox}Do generalized eigenspaces account for every vector? For $T(x)=(4x_2,0,5x_3)$, $\mathbb C^3=G(0,T)\oplus G(5,T)$ even though ordinary eigenspaces do not span.\end{qbox}
\begin{proposition*}[Notes Proposition 73; LADR 8.9, 8.22]For a finite-dimensional complex space and all distinct eigenvalues $\lambda_1,\ldots,\lambda_r$, $V=\bigoplus_jG(\lambda_j,T)$. Each summand is invariant; on it $T=\lambda_jI+N_j$ with $N_j$ nilpotent. In particular, $V$ has a basis of generalized eigenvectors.\end{proposition*}
\begin{motive}\ii[Want.] Prove generalized eigenspaces span, as their directness is already known.\ii[Know.] Complex operators in nonzero finite dimension have eigenvalues; stabilized kernel and range are complementary.\ii[Key step.] Separate one nonzero generalized eigenspace and induct on the proper invariant range.\end{motive}
\begin{prob}\begin{proof}\begin{steps}\ii Induct on $n=\dim V$, with the empty basis at $n=0$. For $n>0$ choose an eigenvalue $\lambda$. Apply LADR 8.4 to $A=T-\lambda I$: $V=G(\lambda,T)\oplus U$, where $U=\im(T-\lambda I)^n$ is invariant.\ii The first summand is nonzero, so $\dim U<n$. Induction supplies a basis of $U$ made of generalized eigenvectors of $T|_U$, also generalized eigenvectors of $T$.\ii Any basis of $G(\lambda,T)$ consists of generalized eigenvectors. Adjoining it gives a generalized eigenvector basis of $V$, proving that the sum of all $G(\lambda_j,T)$ spans. Proposition 72 makes the sum direct.\ii Each summand is invariant by Proposition 70. On it $(T-\lambda_jI)^n=0$, so $N_j=(T-\lambda_jI)|_{G(\lambda_j,T)}$ is nilpotent, as asserted.\end{steps}\end{proof}\end{prob}
\begin{prob}\begin{definition*}[Notes Definition 38; LADR 8.23, 8.26]The algebraic multiplicity of $\lambda$ is $d_\lambda=\dim G(\lambda,T)$; its geometric multiplicity is $g_\lambda=\dim E(\lambda,T)$ (Week 7). Over $\mathbb C$, the characteristic polynomial is $\chi_T(x)=\prod_\lambda(x-\lambda)^{d_\lambda}$ over the distinct eigenvalues. The empty product for $V=0$ is 1.\end{definition*}\end{prob}
\begin{proposition*}[Notes Definition 38, consequences; LADR 8.25, 8.28]Over $\mathbb C$, $\sum_\lambda d_\lambda=\dim V$, $g_\lambda\le d_\lambda$, and $T$ is diagonalizable exactly when $g_\lambda=d_\lambda$ for every eigenvalue. Also $\deg\chi_T=\dim V$, and its roots are exactly the eigenvalues.\end{proposition*}
\begin{motive}\ii[Want.] Relate multiplicities to the existence of an eigenbasis.\ii[Know.] Generalized eigenspaces form a direct decomposition, and ordinary eigenspaces are contained in them.\ii[Key step.] Compare the dimensions of these two direct sums.\end{motive}
\begin{prob}\begin{proof}\begin{steps}\ii Proposition 73 and the dimension formula for a direct sum give $\sum d_\lambda=\dim V$. Containment $E(\lambda,T)\subseteq G(\lambda,T)$ gives $g_\lambda\le d_\lambda$.\ii By Week 7, an eigenbasis exists exactly when $\sum g_\lambda=\dim V=\sum d_\lambda$. Since each difference is nonnegative, this holds exactly when every difference is zero.\ii The defining product for $\chi_T$ has degree $\sum d_\lambda$ and roots precisely the listed eigenvalues, proving the remaining claims.\end{steps}\end{proof}\end{prob}
\begin{proposition*}[Notes Proposition 74; LADR 8.29--8.30: Cayley--Hamilton]Over $\mathbb C$, $\chi_T(T)=0$. Consequently $m_T\mid\chi_T$.\end{proposition*}
\begin{motive}\ii[Want.] Prove that the characteristic polynomial annihilates the whole operator.\ii[Know.] On each generalized eigenspace the shifted operator is nilpotent; its exponent is at most that summand's dimension.\ii[Key step.] Use dimension $d_\lambda$ on the restriction, rather than the larger ambient dimension.\end{motive}
\begin{prob}\begin{proof}\begin{steps}\ii On $G(\lambda,T)$, the nilpotent $N_\lambda$ satisfies $N_\lambda^{d_\lambda}=0$ by LADR 8.16. Hence $(T-\lambda I)^{d_\lambda}$ is zero on that summand.\ii The commuting product defining $\chi_T(T)$ contains this factor, so is zero on each generalized eigenspace. Proposition 73 says their sum is all of $V$, so $\chi_T(T)=0$.\ii Week 7's divisibility criterion now gives $m_T\mid\chi_T$. The case $V=0$ gives $1(T)=I=0$.\end{steps}\end{proof}\end{prob}
\begin{proposition*}[LADR 8.31: multiplicities in triangular form]In any upper-triangular matrix of a complex operator, the number of diagonal appearances of $\lambda$ is $d_\lambda$. Thus $\chi_T$ is the product of the linear factors from its diagonal entries.\end{proposition*}
\begin{motive}\ii[Want.] Compute algebraic multiplicities from any triangular basis.\ii[Know.] The multiplicities sum to dimension, and powers of triangular matrices remain triangular.\ii[Key step.] Bound each kernel dimension by the number of zero diagonal entries, then sum the bounds.\end{motive}
\begin{prob}\begin{proof}\begin{steps}\ii In a triangular matrix with $b$ zero diagonal entries, the columns with nonzero diagonal are independent: in a proposed relation, the greatest selected index with nonzero coefficient has a nonzero coordinate there. Thus rank is at least $n-b$, and nullity at most $b$.\ii Products of triangular matrices are triangular and their diagonal entries multiply, since $(AB)_{ii}=\sum_k a_{ik}b_{ki}$ has only $k=i$ nonzero. Therefore $(T-\lambda I)^n$ has as many zero diagonal entries as $T$ has entries equal to $\lambda$.\ii Step 1 gives $d_\lambda=\dim\ker(T-\lambda I)^n\le b_\lambda$, its diagonal count. Every diagonal entry is an eigenvalue (Week 7), so $\sum b_\lambda=n=\sum d_\lambda$.\ii Each nonnegative difference $b_\lambda-d_\lambda$ must be zero. Substitute these multiplicities into the defining product for $\chi_T$.\end{steps}\end{proof}\end{prob}
\begin{prob}\begin{definition*}[LADR 8.35]A block-diagonal matrix $\operatorname{diag}(A_1,\ldots,A_r)$ has square matrices $A_j$ along its diagonal and zeros outside these blocks. It represents an operator preserving the corresponding direct-sum coordinate subspaces.\end{definition*}\end{prob}
\begin{proposition*}[LADR 8.37]A complex operator has a block-diagonal matrix with one upper-triangular block for each eigenvalue $\lambda$, of size $d_\lambda$, with every diagonal entry in that block equal to $\lambda$.\end{proposition*}
\begin{motive}\ii[Want.] Express the generalized decomposition as matrix blocks.\ii[Know.] $T=\lambda I+N_\lambda$ on each invariant summand, with $N_\lambda$ nilpotent.\ii[Key step.] Choose a strictly triangular basis for each nilpotent restriction and concatenate the bases.\end{motive}
\begin{prob}\begin{proof}\begin{steps}\ii By the nilpotence tests, each $N_\lambda$ has a strictly upper-triangular matrix. Adding $\lambda I$ gives the required diagonal and block size $\dim G(\lambda,T)=d_\lambda$.\ii The generalized eigenspace direct sum makes the concatenation a basis. Invariance of each summand makes every off-block entry zero.\end{steps}\end{proof}\end{prob}
\begin{example*}[LADR 8.24, 8.27, 8.38]\begin{steps}\ii For $A=\left(\begin{smallmatrix}6&3&4\\0&6&2\\0&0&7\end{smallmatrix}\right)$, $G(6,A)=\operatorname{span}(e_1,e_2)$ and $G(7,A)=\mathbb C(10,2,1)$, by the shifted kernel equations.\ii The multiplicities are 2 and 1, so $\chi_A=(x-6)^2(x-7)$. In basis $e_1,e_2,(10,2,1)$ the matrix is $\operatorname{diag}(\left(\begin{smallmatrix}6&3\\0&6\end{smallmatrix}\right),[7])$, checked from the three images.\end{steps}\end{example*}
\begin{example*}[LADR 8.36]The five-by-five matrix $\operatorname{diag}([4],\left(\begin{smallmatrix}2&-3\\0&2\end{smallmatrix}\right),\left(\begin{smallmatrix}1&7\\0&1\end{smallmatrix}\right))$ has three upper-triangular constant-diagonal blocks. Its characteristic polynomial is $(x-4)(x-2)^2(x-1)^2$ by the diagonal-count result.\end{example*}
\begin{shortcuts}\ii For a triangular complex matrix, count diagonal occurrences to obtain $\chi_T$; no determinant expansion is required.\ii Algebraic multiplicity is the stabilized shifted-kernel dimension; geometric multiplicity uses only the first kernel.\end{shortcuts}
\begin{trapbox}\ii Algebraic and geometric multiplicities can differ: $\left(\begin{smallmatrix}1&1\\0&1\end{smallmatrix}\right)$ has algebraic multiplicity 2 but geometric multiplicity 1.\ii Cayley--Hamilton does not make every vector an eigenvector: for $T=\operatorname{diag}(1,2)$, $\chi_T(T)(1,1)=0$, while $T(1,1)=(1,2)$ is not a scalar multiple of $(1,1)$.\end{trapbox}
\section{Jordan blocks and their sizes}
\begin{qbox}How do the shifted kernels specify a canonical matrix? A three-vector chain for $T-I$ gives $\left(\begin{smallmatrix}1&1&0\\0&1&1\\0&0&1\end{smallmatrix}\right)$, with successive kernel dimensions $1,2,3$.\end{qbox}
\begin{prob}\begin{definition*}[Notes Definition 40; LADR 8.44]The Jordan block $J_r(\lambda)$ is the $r$-by-$r$ matrix with $\lambda$ on its diagonal, 1 immediately above the diagonal and zeros elsewhere. A Jordan normal form is a block-diagonal matrix of such blocks; a basis giving that matrix is a Jordan basis. Different blocks may have the same eigenvalue. A block is $\lambda I+N$ with one length-$r$ chain for $N$.\end{definition*}\end{prob}
\begin{proposition*}[Notes Proposition 76; LADR 8.46]Every operator on a finite-dimensional complex space has a Jordan basis.\end{proposition*}
\begin{motive}\ii[Want.] Improve upper-triangular blocks to unit-chain blocks.\ii[Know.] Generalized eigenspaces form a direct sum and each shifted restriction is nilpotent.\ii[Key step.] Apply the nilpotent chain-basis theorem separately on each summand.\end{motive}
\begin{prob}\begin{proof}\begin{steps}\ii On $G(\lambda,T)$, choose a chain basis for $N_\lambda=(T-\lambda I)|_G$ by Proposition 75. Its matrix consists of zero-diagonal unit-chain blocks.\ii Adding $\lambda I$ turns each into $J_r(\lambda)$. Concatenating bases over the direct sum in Proposition 73 gives a Jordan basis of $V$.\end{steps}\end{proof}\end{prob}
\begin{example*}[Notes Chapter 15: displayed Jordan forms]The three five-by-five matrices displayed after Definition 40 are $\operatorname{diag}(J_2(2),J_3(3))$, $\operatorname{diag}(J_2(2),J_2(3),J_1(3))$, and $\operatorname{diag}(2,2,3,3,3)$. They have respectively two, three and five Jordan blocks; the last is diagonal.\end{example*}
\begin{proposition*}[Jordan data: minimal polynomial and multiplicities]For each eigenvalue $\lambda$, let its block lengths be $r_{\lambda,1},\ldots,r_{\lambda,b_\lambda}$. Its geometric multiplicity is $b_\lambda$, its algebraic multiplicity is $\sum_jr_{\lambda,j}$, and its factor in $m_T$ is $(x-\lambda)^{\max_jr_{\lambda,j}}$. Thus $T$ is diagonalizable exactly when every block has size 1.\end{proposition*}
\begin{motive}\ii[Want.] Read the two multiplicities and minimal polynomial directly from blocks.\ii[Know.] The generalized eigenspace consists of its blocks, and each shifted block has a one-dimensional kernel.\ii[Key step.] A polynomial annihilates a size-$r$ block exactly when its root at $\lambda$ has multiplicity at least $r$.\end{motive}
\begin{prob}\begin{proof}\begin{steps}\ii A size-$r$ block has $\ker(J_r(\lambda)-\lambda I)=Fe_1$ and its full $r$-dimensional space is generalized. On other-eigenvalue blocks the shift is invertible, since their triangular diagonal entries are nonzero. This proves both multiplicity formulas.\ii Write $p=(x-\lambda)^sq$ with $q(\lambda)\ne0$. Polynomial division gives $q(\lambda I+N)=q(\lambda)I+NR(N)$ for a polynomial $R$.\ii This last operator is invertible: after factoring out $q(\lambda)$, it is $I+K$ with $K$ nilpotent; $(I+K)\sum_{j=0}^{r-1}(-K)^j=I-(-K)^r=I$. Thus $p(J)=N^sq(J)$ is zero precisely when $s\ge r$.\ii A polynomial annihilates a block-diagonal matrix exactly when it annihilates every block. The least monic product with the required multiplicity for each root has the stated maximal exponents. Apply Week 7's criterion for the final claim.\end{steps}\end{proof}\end{prob}
\begin{proposition*}[Notes Proposition 77]For integers $k\ge1$, put $d_k=\dim\ker(T-\lambda I)^k$, with $d_0=0$. The number of blocks for $\lambda$ of size at least $k$ is $d_k-d_{k-1}$; the number of size exactly $k$ is $2d_k-d_{k-1}-d_{k+1}$. Consequently Jordan form is unique up to block order.\end{proposition*}
\begin{motive}\ii[Want.] Recover all block sizes without constructing a basis.\ii[Know.] The $k$th power of a length-$r$ shifted block has nullity $\min(k,r)$.\ii[Key step.] First differences count blocks still contributing a new kernel direction; second differences count those that stop.\end{motive}
\begin{prob}\begin{proof}\begin{steps}\ii On a block for $\lambda$, the shift sends $e_j$ to $e_{j-1}$ and the first $\min(k,r)$ vectors to zero after $k$ steps. Thus its nullity is $\min(k,r)$. Other-eigenvalue blocks have invertible shift and contribute zero.\ii Subtracting the nullities at $k$ and $k-1$ contributes 1 for each block with $r\ge k$, and 0 for each smaller block. Summing proves the first formula.\ii Subtract the counts for sizes at least $k+1$ from sizes at least $k$: $(d_k-d_{k-1})-(d_{k+1}-d_k)=2d_k-d_{k-1}-d_{k+1}$.\ii These numbers determine every block multiplicity for every eigenvalue, proving uniqueness up to order.\end{steps}\end{proof}\end{prob}
\method{Construct an actual Jordan basis (Notes Example 62)}\begin{steps}\ii For $A=\left(\begin{smallmatrix}1&2&3\\0&1&2\\0&0&1\end{smallmatrix}\right)$ set $N=A-I$. Then $N^2=\left(\begin{smallmatrix}0&0&4\\0&0&0\\0&0&0\end{smallmatrix}\right)$ and $N^3=0$. The successive nullities are $1,2,3$, so one length-3 block occurs.\ii Choose $v_3=e_3$ outside $\ker N^2$. Then $v_2=Nv_3=(3,2,0)$ and $v_1=Nv_2=(4,0,0)$.\ii With $S=(v_1\ v_2\ v_3)=\left(\begin{smallmatrix}4&3&0\\0&2&0\\0&0&1\end{smallmatrix}\right)$, the chain equations give $AS=SJ_3(1)$, hence $A=SJ_3(1)S^{-1}$.\end{steps}
\method{Recover sizes from nullities (Notes Example 63)}\begin{steps}\ii On $\mathbb C^8$, the shifted nullities for eigenvalue 1 are $d_0,d_1,d_2,d_3=0,2,4,4$. The differences are $2,2,0$, so both blocks have length 2.\ii For eigenvalue 2 they start $0,2,3,4$. These account for dimension 4, since the first generalized eigenspace already has dimension 4. The next nullity is therefore 4, and differences are $2,1,1,0$. The sizes are 3 and 1.\ii Thus Jordan form is $\operatorname{diag}(J_2(1),J_2(1),J_3(2),J_1(2))$. The characteristic polynomial is $(x-1)^4(x-2)^4$ and the minimal polynomial is $(x-1)^2(x-2)^3$.\end{steps}
\begin{example*}[Notes Chapter 15: generalized bases do not determine similarity]\begin{steps}\ii For $A=\left(\begin{smallmatrix}3&-1&-1\\2&0&-1\\2&-1&0\end{smallmatrix}\right)$, $N=A-I$ has rank 1, $N^2=0$ and $N\ne0$. Thus its blocks have lengths 2 and 1.\ii For $B=\left(\begin{smallmatrix}3&-2&0\\1&1&-1\\0&2&-1\end{smallmatrix}\right)$, $M=B-I$ has rank 2, $M^2\ne0$ and $M^3=0$, so its sole block has length 3.\ii For both matrices every standard basis vector is a generalized eigenvector for 1, as it is for $I_3$. Nevertheless their shifted ranks are respectively 1, 2 and 0. Similarity preserves rank because an invertible change of coordinates preserves range dimension (Week 4), so these three matrices are pairwise nonsimilar.\end{steps}\end{example*}
\begin{shortcuts}\ii Kernel first differences must be nonnegative and nonincreasing: they count blocks of size at least $k$. This quickly detects impossible data.\ii The first nullity counts blocks, the final nullity gives their total size, and the last increase locates the largest size.\ii To construct a longest chain, choose a vector in $\ker(T-\lambda I)^r$ outside the preceding kernel and apply the shift repeatedly.\end{shortcuts}
\begin{trapbox}\ii Algebraic multiplicity need not equal the exponent in the minimal polynomial: two blocks $J_2(1),J_2(1)$ give algebraic multiplicity 4 but exponent 2.\ii A generalized eigenvector basis need not be a Jordan basis: the standard basis for the matrix in Example 62 consists of generalized eigenvectors, but its nonunit entries do not have Jordan form.\end{trapbox}
\section{Powers and square roots}
\begin{qbox}How do Jordan blocks make repeated application explicit? For $J=\left(\begin{smallmatrix}1&1&0\\0&1&1\\0&0&1\end{smallmatrix}\right)$, $J^2=\left(\begin{smallmatrix}1&2&1\\0&1&2\\0&0&1\end{smallmatrix}\right)$; the second superdiagonal records a quadratic factor in the exponent.\end{qbox}
\begin{proposition*}[Notes Chapter 15: powers and Jordan--Chevalley decomposition]For $J_r(\lambda)=\lambda I+N$ and integer $k\ge0$, $J_r(\lambda)^k=\sum_{j=0}^{\min(k,r-1)}\binom{k}{j}\lambda^{k-j}N^j$, with the convention $\lambda^0=1$. A Jordan matrix is $D+N$ with diagonal $D$, nilpotent $N$ and $DN=ND$. If invertible it also equals $D(I+D^{-1}N)$, with commuting factors.\end{proposition*}
\begin{motive}\ii[Want.] Replace large matrix powers with finitely many scalar expressions.\ii[Know.] The diagonal scalar part commutes with the shifted block, whose $r$th power is zero.\ii[Key step.] Apply the binomial theorem for commuting operators and omit the zero powers.\end{motive}
\begin{prob}\begin{proof}\begin{steps}\ii For commuting $A,B$, induction on $k$ gives $(A+B)^k=\sum_{j=0}^k\binom{k}{j}A^{k-j}B^j$: multiply by $A+B$, commute the factors, and combine adjacent terms using Pascal's identity.\ii Take $A=\lambda I$, $B=N$ and use $N^j=0$ for $j\ge r$. Restricting to $j\le k$ avoids negative exponents, including when $\lambda=0$.\ii For a full Jordan matrix, let $D$ equal its diagonal and $N=J-D$. On each block $D$ is a scalar matrix, so it commutes with $N$; powers of the block shifts show $N$ nilpotent.\ii If $J$ is invertible, none of its diagonal eigenvalues is zero, so $D$ is invertible and $D(I+D^{-1}N)=D+N$. Both factors commute because $D$ and $N$ do. These are the additive and multiplicative decompositions described in the notes.\end{steps}\end{proof}\end{prob}
\method{Compute the hundredth power (Notes Chapter 15)}\begin{steps}\ii For $J=J_3(1)$, $N=J-I$ has $N^3=0$, so $J^{100}=I+100N+\binom{100}{2}N^2$.\ii Since $\binom{100}{2}=100\cdot99/2=4950$, $J^{100}=\left(\begin{smallmatrix}1&100&4950\\0&1&100\\0&0&1\end{smallmatrix}\right)$.\end{steps}
\begin{prob}\begin{definition*}[LADR 8C: square root]An operator $R$ is a square root of $T$ if $R^2=T$.\end{definition*}\end{prob}
\begin{proposition*}[LADR 8.39]Over $\mathbb R$ or $\mathbb C$, if $N$ is nilpotent, then $I+N$ has a square root which is a polynomial in $N$.\end{proposition*}
\begin{motive}\ii[Want.] Solve an operator square equation using the known nilpotent powers.\ii[Know.] High powers of $N$ vanish, so only finitely many polynomial coefficients matter.\ii[Key step.] Choose each coefficient to cancel the next unwanted power in the square.\end{motive}
\begin{prob}\begin{proof}\begin{steps}\ii Choose $r$ with $N^r=0$. If $r=1$, use $R=I$. Otherwise seek $R=I+\sum_{j=1}^{r-1}a_jN^j$.\ii The coefficient of $N$ in $R^2$ is $2a_1$, so set $a_1=1/2$. For $2\le k<r$, set $a_k=-\frac12\sum_{j=1}^{k-1}a_ja_{k-j}$.\ii Then the coefficient of $N^k$ in $R^2$ is $2a_k+\sum_{j=1}^{k-1}a_ja_{k-j}=0$. All terms of degree at least $r$ vanish. Hence $R^2=I+N$. The first coefficients are $1/2,-1/8,1/16$.\end{steps}\end{proof}\end{prob}
\begin{proposition*}[LADR 8.41]Every invertible operator on a finite-dimensional complex vector space has a square root.\end{proposition*}
\begin{motive}\ii[Want.] Reduce a general square-root problem to identity plus nilpotent.\ii[Know.] On $G(\lambda,T)$, $T=\lambda I+N_\lambda$, and invertibility gives $\lambda\ne0$. Every complex scalar has a square root (Week 6).\ii[Key step.] Factor out $\lambda$, construct a root on each summand, then combine them by the direct decomposition.\end{motive}
\begin{prob}\begin{proof}\begin{steps}\ii On $G(\lambda,T)$ write $T=\lambda(I+N_\lambda/\lambda)$. The scaled operator $N_\lambda/\lambda$ is nilpotent, so the preceding result gives $S_\lambda^2=I+N_\lambda/\lambda$.\ii Choose $c_\lambda\in\mathbb C$ with $c_\lambda^2=\lambda$. Set $R_\lambda=c_\lambda S_\lambda$, so $R_\lambda^2=T|_{G(\lambda,T)}$.\ii Define $R(\sum_\lambda v_\lambda)=\sum_\lambda R_\lambda v_\lambda$ using Proposition 73's unique decomposition. This is linear and $R^2v=\sum_\lambda Tv_\lambda=Tv$.\end{steps}\end{proof}\end{prob}
\method{Construct a square root}\begin{steps}\ii For $J_3(1)=I+N$, $N^3=0$, so $R=I+N/2-N^2/8=\left(\begin{smallmatrix}1&1/2&-1/8\\0&1&1/2\\0&0&1\end{smallmatrix}\right)$.\ii Expanding gives $R^2=I+N+(1/4-1/4)N^2=J_3(1)$, since all higher powers vanish.\end{steps}
\begin{shortcuts}\ii Compute a power block by block, using only powers of each nilpotent shift up to its block size minus 1.\ii A square root of $I+N$ is found by matching coefficients; no infinite-series convergence is needed because the operator expression is finite.\end{shortcuts}
\begin{trapbox}\ii An invertible real operator need not have a real square root: multiplication by $-1$ on $\mathbb R$ would require a real scalar $r$ with $r^2=-1$.\ii A complex noninvertible operator may have no square root: if $R^2=J_3(0)$, then $R^6=0$, so the dimension bound gives $R^3=0$. This forces $J_3(0)^2=R^4=0$, contradicting its nonzero second power.\end{trapbox}
\section{Problem set}
\begin{qbox}Recover blocks from shifted kernels and use them for powers and minimal polynomials. For example, nullities $0,1,2,3,3$ give one block of length 3.\end{qbox}
\begin{prob}\textbf{HW8 P1.} For differentiation $D:\mathcal P_n(\mathbb C)\to\mathcal P_n(\mathbb C)$, (a) prove nilpotence and determine its index; (b) determine the geometric and algebraic multiplicities of 0; (c) find a basis giving a single Jordan block.\\\emph{Hint:} Use the derivative chain $1,x,x^2/2!,\ldots$ and the nilpotence tests.\end{prob}
\begin{prob}\textbf{HW8 P2.} Construct explicit matrices with (a) eigenvalue 2 of algebraic multiplicity 3 and geometric multiplicity 1, and eigenvalue $-1$ of algebraic and geometric multiplicity 2; (b) the same algebraic multiplicities but geometric multiplicity of 2 equal to 2. Give each minimal polynomial.\\\emph{Hint:} Use the Jordan-data formulas for number and largest size of blocks.\end{prob}
\begin{prob}\textbf{HW8 P3.} For $J=\left(\begin{smallmatrix}1&1&0\\0&1&1\\0&0&1\end{smallmatrix}\right)$, (a) compute $\ker(J-I)$, $\ker(J-I)^2$, $\ker(J-I)^3$; (b) determine $G(1,J)$ and both multiplicities; (c) compute $J^m$ for each positive integer $m$.\\\emph{Hint:} Use the chain basis and the commuting binomial formula.\end{prob}
\begin{prob}\textbf{HW8 P4.} A nilpotent operator has one block of size 3 and one of size 2. (a) Compute $\dim\ker N^k$ for $k=0,1,2,3,4$. (b) On the same five-dimensional space, if instead given only $\dim\ker N=2$, $\dim\ker N^2=4$, $\dim\ker N^3=5$, explain how these force the same two block sizes.\\\emph{Hint:} Apply Proposition 77's first differences to these five-dimensional data.\end{prob}
\begin{prob}\textbf{HW8 P5 (optional).} On $\mathbb C^6$, $T$ has eigenvalues 1 and 2 with shifted nullities $1,2,2$ at powers $1,2,3$ for 1, and $2,3,4,4$ at powers $1,2,3,4$ for 2. Determine its Jordan form up to block order.\\\emph{Hint:} Stabilization and Proposition 77 recover lengths from successive differences.\end{prob}
\begin{prob}\textbf{HW8 P6 (optional).} For blocks $J_3(4),J_2(4),J_1(-1)$, determine the characteristic polynomial, minimal polynomial, both multiplicities of each eigenvalue, and diagonalizability.\\\emph{Hint:} Use sums of sizes for $\chi_T$ and maxima of sizes for $m_T$.\end{prob}
\begin{prob}\textbf{HW8 P7 (optional).} Prove that for nilpotent $N$, $\dim\ker N^k-\dim\ker N^{k-1}$ counts the Jordan blocks of size at least $k$.\\\emph{Hint:} Use the single-block nullity $\min(k,r)$ from Proposition 77.\end{prob}
\begin{prob}\textbf{Exam practice (midterm2-review.pdf, Problem 5).} For $A=\left(\begin{smallmatrix}3&1&1\\0&3&-1\\0&0&1\end{smallmatrix}\right)$ as an operator $\mathbb R^3\to\mathbb R^3$, find its minimal polynomial, characteristic polynomial, eigenvalues and eigenspaces. For the characteristic polynomial, regard the same real matrix on $\mathbb C^3$ and use the intrinsic definition above.\\\emph{Hint:} Use triangular diagonal counts, then solve shifted kernels and test the candidate annihilating factors.\end{prob}
\end{document}
